%===========================================================================
\chapter{Discussion}\label{ch:disc}
%===========================================================================

\section{Design trade-offs}
\textbf{Dedicated accumulator vs.\ register-file accumulation.} A
dedicated 64-bit accumulator gives full precision and needs only two
register-file read ports. The cost is extra architectural state, which
an operating system must save on a context switch, and explicit
\code{rdacc} instructions to read the result. The Xpulp
\code{p.mac}~\cite{cv32e40p} approach needs a third read port and
accumulates in only 32 bits.

\textbf{Saturation vs.\ wrap-around.} Signed operations saturate,
which is the expected DSP behaviour, whereas \code{macu} wraps modulo $2^{64}$,
which is what checksums and hashing need.

\textbf{Two-stage MAC pipeline aligned with the CPU pipeline.} Putting
the multiplication in EX and the accumulation in MEM means MAC
instructions issue back-to-back with no forwarding inside the unit.
Only a consumer placed \emph{directly} after a MAC result pays one
stall cycle, exactly like a load, and in the scheduled FIR code this never
happens.

\textbf{Instruction count vs.\ loop overhead.} The MAC instructions
cut the arithmetic work per tap from six instructions to one. In the
looped FIR, however, loads and loop control remain, so the measured
speed-up is 1.59$\times$. Once the loop overhead is removed by
unrolling, the MAC unit's advantage grows to 2.67$\times$. Adding
post-increment loads or hardware loops (as in Xpulp) would bring the
looped code close to the unrolled performance.

\section{Limitations}
The synthesis results come from open-source tools and a generic gate library.
No commercial ASIC standard-cell library was used, so the report gives gate
counts and logic depth rather than area in $\mu$m$^2$ or power figures. The MAC
unit's critical path is the 34$\times$34-bit combinational Booth multiplier
(17.7\,ns on ECP5, 127 gate levels), whose partial products are summed with
ripple-carry adders. The processor implements the RV32M multiply instructions
in a single cycle, which is simple but sets the processor's critical path
(28.9\,ns, 34.6\,MHz on ECP5). It has no caches, interrupts, CSRs or division
instructions, and its memories are modelled with combinational reads; on an
FPGA they would be replaced by synchronous block RAM, which requires an extra
fetch stage or registered addresses. Finally, the programs are written in
assembly with a project-specific assembler, and C code would need compiler
support for the custom instructions, for example through inline assembly.

\section{Future work}
Several extensions would build on this work. The critical paths could be
shortened by using a carry-save (Wallace/Dadda) tree with a parallel-prefix
final adder (for example Kogge-Stone) for the Booth multiplier, and by
pipelining the RV32M multiplier over EX and MEM, as is already done for the MAC
unit. Post-increment loads and a zero-overhead hardware loop could be added so
that looped DSP code reaches the throughput of the unrolled code, and
packed-SIMD operations (for example two 16$\times$16 MACs per cycle), similar
to the RISC-V P-extension proposals, would widen the arithmetic. The
instruction and data memories could be mapped to FPGA block RAM so that the
processor runs on a real ECP5 board, and the design could be taken through an
ASIC flow with a real standard-cell library (for example OpenLane with the
SkyWater 130\,nm PDK). Finally, the Booth multiplier could be formally verified
against \code{*}, for example with SymbiYosys.

%===========================================================================
\chapter{Conclusion}\label{ch:concl}
%===========================================================================
This project designed a RISC-V custom-extension multiply-accumulate unit for
digital signal processing, integrated it into a complete 5-stage pipelined
RISC-V processor, and verified both thoroughly. The MAC unit adds eight
\code{custom-0} instructions and is built around a radix-4 Booth multiplier
and a 64-bit saturating accumulator with Q15 support. Its two-stage pipeline
accepts one instruction per cycle and lines up exactly with the EX and MEM
stages of the processor.

At the unit level, an independent reference model and scoreboard, exhaustive
multiplier tests, directed corner cases, three analytical DSP benchmarks and a
3{,}000-instruction constrained-random regression performed 88{,}896 checks,
all with zero errors.

At the system level, real RISC-V programs, built with a project-specific
assembler, ran on the processor: an 18-part ISA and hazard self-test, and a
16-tap low-pass FIR over 128 samples in four versions. Every program passed.
The FIR outputs were bit-identical to the golden model, the unwanted tone was
attenuated by 60.1\,dB as the filter design predicts, and mutation tests
confirmed that the system tests detect forwarding and stall bugs. Compared with
the same filter in plain RV32IM software, the MAC instructions made the FIR
1.59$\times$ faster in looped form and 2.67$\times$ faster in unrolled form,
with 2.6$\times$ smaller code.

Yosys synthesis produced a latch-free MAC netlist of 9{,}532 gates and 179
flip-flops, which passes the full testbench at gate level. On a Lattice ECP5
FPGA, the MAC unit reaches 56.4\,MHz, and the complete processor with the
MAC unit reaches 34.6\,MHz using 9\,\% of the device's logic. The schematics,
floorplans and timing analysis show that the multipliers, both the Booth
multiplier and the RV32M multiplier, are the largest blocks and set the
critical paths. That makes them the main target for future optimisation.
